\documentclass[11pt,a4paper]{article}
\usepackage[utf8]{inputenc}
\usepackage[portuguese]{babel}
\usepackage{geometry}
\usepackage{fancyhdr}
\usepackage{hyperref}
\usepackage{listings}
\usepackage{xcolor}

\geometry{left=2cm,right=2cm,top=2cm,bottom=2cm}
\setlength{\parskip}{0.4em}
\pagestyle{fancy}
\fancyhf{}
\lhead{CC8550 --- Simulação e Teste de Software}
\rhead{Laboratório 06}
\cfoot{\thepage}

\lstset{
    basicstyle=\ttfamily\small,
    breaklines=true,
    frame=single,
    backgroundcolor=\color{gray!10}
}

\title{
    {\Large \textbf{Centro Universitário FEI}}\\
    \vspace{0.3cm}
    {\large Ciência da Computação}\\[0.8cm]
    {\LARGE \textbf{Simulação e Teste de Software}}\\[0.3cm]
    {\Large Atividade 06 --- Testes de Unidade e Integração}\\[0.8cm]
    {\large \textbf{Calculadora com Repositório de Histórico}}
}
\author{
    \textbf{Aluno:} Lucas Rebouças Silva\\
    \textbf{R.A.:} 22.122.048-6\\[0.4cm]
    \textbf{Professor:} Luciano Rossi
}
\date{2026}

\begin{document}
\maketitle
\thispagestyle{empty}
\vspace{1cm}
\begin{center}
\small
Documento em estilo ABNT simplificado para entrega acadêmica. Conteúdo alinhado ao relatório em \texttt{relatorio.md}.
\end{center}
\newpage
\setcounter{page}{1}

\section{Objetivo}
Desenvolver testes de unidade e de integração para o sistema composto por \texttt{Calculadora} e \texttt{HistoricoRepositorio}, utilizando \emph{test doubles} (stub e mock) conforme a disciplina CC8550.

\section{Bug corrigido em \texttt{potencia}}
No código original do enunciado, o histórico registrava a operação com o operador \texttt{*} em vez de \texttt{**}, gerando strings semanticamente incorretas (ex.: \texttt{2 * 8 = 256}). A correção consiste em persistir:
\begin{lstlisting}[language=Python]
self.repositorio.salvar(f"{base} ** {expoente} = {resultado}")
\end{lstlisting}

\section{Resultados dos testes}
Execução: \texttt{python -m unittest discover tests -v}. Todos os testes passaram (59 testes).

\section{Cobertura em \texttt{calculadora.py}}
Comando: \texttt{python -m coverage run -m unittest discover tests} seguido de \texttt{python -m coverage report -m --include=src/calculadora.py}.

\subsection{Print simulado (100\% de cobertura)}
\begin{lstlisting}
Name                 Stmts   Miss  Cover   Missing
--------------------------------------------------
src\calculadora.py      43      0   100%
--------------------------------------------------
TOTAL                   43      0   100%
\end{lstlisting}

Nenhuma linha permaneceu sem cobertura. Justificativa: todos os ramos de validação de tipo, divisão por zero e operações com registro no repositório foram exercitados.

\section{Reflexão: Stub vs Mock}
\begin{itemize}
    \item \textbf{Stub:} substitui o repositório para isolar a \texttt{Calculadora}; prioriza o resultado numérico sem depender de persistência real.
    \item \textbf{Mock:} verifica interações --- por exemplo, se \texttt{salvar} foi chamado com a string exata esperada e se \emph{não} foi chamado quando uma exceção ocorre antes do registro.
\end{itemize}
Na prática, mocks foram decisivos para revelar o defeito textual na potência, pois o valor de \texttt{base ** expoente} estava correto, mas a string armazenada não.

\end{document}
